% Taller 3 -- Punto 3: revisión (análisis) del script del profesor.   Compilar:  make punto3   (tectonic, plantilla UD)
\documentclass[11pt,letterpaper]{article}
\usepackage{ud-informe}
% Acento azul: el logo del programa (MCIC) es azul, así que se usa la gama azur del manual de marca en lugar del rojo.
% Para volver al rojo institucional de la plantilla, borrar estas dos líneas.
\colorlet{udrojo}{udazul}
\definecolor{udrojoosc}{HTML}{0C6F96}


\renewcommand{\UDtitulo}{Revisión del script del profesor}
\renewcommand{\UDsubtitulo}{Análisis de la solución de referencia: MLP y CNN en PyTorch sobre MNIST}
\renewcommand{\UDtipo}{Taller 3 -- Punto 3 -- Informe de ejecución}
\renewcommand{\UDasignatura}{Inteligencia Computacional Aplicada}
\renewcommand{\UDdocente}{Cesar Andrey Perdomo Charry}
\renewcommand{\UDestudiantes}{Juan Felipe Rodríguez Galindo -- 20261595004}
\renewcommand{\UDencabezado}{Taller 3 -- Punto 3: revisión del script del profesor}

\begin{document}
\UDportada

%% ---------------------------------------------------------------------------------------------
\section{Qué es el script}

El script es el cuaderno \texttt{Taller - DeepLearning PYTHON.ipynb}, escrito en Python con PyTorch y torchvision y ejecutado en Google Colab con GPU.
Entrena y compara un perceptrón multicapa (MLP) y una red convolucional (CNN) sobre \textbf{MNIST} (70\,000 imágenes de dígitos de $28\times28$, 10 clases) y recorre el
flujo completo de un proyecto de clasificación de imágenes: datos, línea base, modelo convolucional, evaluación honesta y un experimento sobre la tasa de aprendizaje.
Esta revisión describe qué hace cada bloque, qué idea enseña y qué límites tiene; las cifras son las que el propio cuaderno dejó guardadas en sus salidas.


%% ---------------------------------------------------------------------------------------------
\section{Diagrama de secuencia}

La figura~\ref{fig:secuencia} muestra quién habla con quién durante la ejecución. El \textbf{cuaderno} es el orquestador: pide los datos a \texttt{torchvision}, crea los lotes con el
\texttt{DataLoader} y, dentro del bucle de entrenamiento, repite los cinco pasos de cada iteración (borrar gradientes con \texttt{Adam}, propagar hacia adelante en el \texttt{Modelo},
calcular la pérdida, retropropagar y actualizar). El conjunto de prueba aparece una sola vez, después de los bucles.

\begin{figure}[H]
\centering
% Diagrama de secuencia (UML) del cuaderno resuelto del profesor. Se incluye con % Diagrama de secuencia (UML) del cuaderno resuelto del profesor. Se incluye con % Diagrama de secuencia (UML) del cuaderno resuelto del profesor. Se incluye con \input{diagrama_secuencia_profesor}
% dentro de un entorno figure. Usa solo la paleta de ud-informe (udsable, udplata, udrojo, udazul).
\begin{tikzpicture}[yscale=0.78, font=\fontsize{7.6pt}{9pt}\selectfont,
    part/.style={draw=udsable!70, fill=white, rounded corners=2pt, minimum width=2.25cm, minimum height=0.7cm, align=center, inner sep=2pt},
    parthecho/.style={part, draw=udrojo, fill=udrojo!8, font=\fontsize{7.6pt}{9pt}\selectfont\bfseries},
    vida/.style={udplata!80, dash pattern=on 2.5pt off 2.5pt, line width=0.5pt},
    sync/.style={-{Latex[length=1.8mm]}, line width=0.65pt, color=udsable},
    retorno/.style={-{Latex[length=1.8mm, open]}, line width=0.55pt, color=udsable!65, dash pattern=on 3pt off 2pt},
    txt/.style={fill=white, inner sep=1.2pt, text=udsable},
    marco/.style={draw=udrojo!80, line width=0.7pt, rounded corners=1pt},
    etiq/.style={fill=udrojoosc, text=white, font=\fontsize{6.8pt}{8pt}\selectfont\bfseries, inner sep=2pt, anchor=north west},
    nota/.style={draw=udplata, fill=udfondo, font=\fontsize{7pt}{8.4pt}\selectfont, align=left, inner sep=3pt}]

  % --- participantes (x en cm) ---
  \def\xU{0.0}   \def\xC{2.95}  \def\xT{5.9}   \def\xD{8.85}  \def\xM{11.8}  \def\xA{14.75}
  \node[part]      (U) at (\xU,0) {\textbf{Usuario}\\(estudiante)};
  \node[parthecho] (C) at (\xC,0) {Cuaderno\\(Jupyter / Colab)};
  \node[part]      (T) at (\xT,0) {\texttt{torchvision}\\(datos)};
  \node[part]      (D) at (\xD,0) {\texttt{DataLoader}\\(lotes)};
  \node[part]      (M) at (\xM,0) {\texttt{Modelo}\\(\texttt{nn.Module})};
  \node[part]      (A) at (\xA,0) {\texttt{Adam}\\(optimizador)};
  \foreach \p/\x in {U/\xU,C/\xC,T/\xT,D/\xD,M/\xM,A/\xA}
    \draw[vida] (\x,-0.35) -- (\x,-16.05);

  % --- marcos de repetición (se dibujan primero para quedar al fondo) ---
  \draw[marco] (\xC-0.45,-5.75) rectangle (\xA+0.55,-13.15);
  \node[etiq] at (\xC-0.45,-5.75) {loop \;\;para cada época (3)};
  \draw[marco, dashed] (\xC-0.25,-6.3) rectangle (\xA+0.35,-11.70);
  \node[etiq] at (\xC-0.25,-6.3) {loop \;\;para cada lote (64 imágenes)};

  % --- mensajes (etiqueta sobre la flecha, sin tapar la línea) ---
  \tikzset{et/.style={txt, above=1.6pt}}
  \draw[sync]    (\xU,-1.5)  -- node[et]{ejecutar celdas} (\xC,-1.5);
  \draw[sync]    (\xC,-2.2)  -- node[et]{\texttt{cargar\_dataset("MNIST")}} (\xT,-2.2);
  \draw[retorno] (\xT,-2.8)  -- node[et]{tensores $N{\times}1{\times}28{\times}28$} (\xC,-2.8);
  \draw[sync]    (\xC,-3.4)  -- (\xC+0.55,-3.4) -- (\xC+0.55,-3.85) -- (\xC,-3.85);
  \node[txt, anchor=west] at (\xC+0.65,-3.62) {\texttt{preparar\_datos()}: 55\,000 / 5\,000 / 10\,000; $x'=(x-\mu)/\sigma$};
  \draw[sync]    (\xC,-4.55) -- node[et]{\texttt{TensorDataset}, lotes de 64} (\xD,-4.55);
  \draw[sync]    (\xC,-5.2)  -- node[et]{\texttt{Adam(parámetros, lr=1e-3)}} (\xA,-5.2);

  \draw[retorno] (\xD,-7.40)  -- node[et]{\texttt{xb, yb}} (\xC,-7.40);
  \draw[sync]    (\xC,-8.00)  -- node[et]{\texttt{zero\_grad()}} (\xA,-8.00);
  \draw[sync]    (\xC,-8.60)  -- node[et]{\texttt{modelo(xb)}} (\xM,-8.60);
  \draw[retorno] (\xM,-9.15) -- node[et]{\emph{logits}} (\xC,-9.15);
  \draw[sync]    (\xC,-9.70)  -- (\xC+0.55,-9.70) -- (\xC+0.55,-10.05) -- (\xC,-10.05);
  \node[txt, anchor=west] at (\xC+0.65,-9.87) {\texttt{F.cross\_entropy(logits, yb)} $\to$ pérdida};
  \draw[sync]    (\xC,-10.80) -- node[et]{\texttt{pérdida.backward()}} (\xM,-10.80);
  \draw[sync]    (\xC,-11.35) -- node[et]{\texttt{opt.step()}} (\xA,-11.35);

  \draw[sync]    (\xC,-12.25) -- node[et]{\texttt{evaluar(validación)}} (\xM,-12.25);
  \draw[retorno] (\xM,-12.80) -- node[et]{pérdida y exactitud} (\xC,-12.80);

  \draw[sync]    (\xC,-13.80) -- node[et]{\texttt{evaluar(prueba)}, \emph{una sola vez}} (\xM,-13.80);
  \draw[retorno] (\xM,-14.35) -- node[et]{matriz de confusión; pares} (\xC,-14.35);
  \draw[sync]    (\xC,-14.95) -- node[et]{\texttt{torch.save} y recargar} (\xM,-14.95);
  \draw[retorno] (\xC,-15.60) -- node[et]{\texttt{[OK]} / \texttt{[REVISAR]}} (\xU,-15.60);

  \node[nota, anchor=north east] at (\xA+0.55,-16.50) {Paso 5: repite el bucle con $\eta\in\{10^{-1},10^{-3},10^{-5}\}$, 2 épocas cada una.};
\end{tikzpicture}

% dentro de un entorno figure. Usa solo la paleta de ud-informe (udsable, udplata, udrojo, udazul).
\begin{tikzpicture}[yscale=0.78, font=\fontsize{7.6pt}{9pt}\selectfont,
    part/.style={draw=udsable!70, fill=white, rounded corners=2pt, minimum width=2.25cm, minimum height=0.7cm, align=center, inner sep=2pt},
    parthecho/.style={part, draw=udrojo, fill=udrojo!8, font=\fontsize{7.6pt}{9pt}\selectfont\bfseries},
    vida/.style={udplata!80, dash pattern=on 2.5pt off 2.5pt, line width=0.5pt},
    sync/.style={-{Latex[length=1.8mm]}, line width=0.65pt, color=udsable},
    retorno/.style={-{Latex[length=1.8mm, open]}, line width=0.55pt, color=udsable!65, dash pattern=on 3pt off 2pt},
    txt/.style={fill=white, inner sep=1.2pt, text=udsable},
    marco/.style={draw=udrojo!80, line width=0.7pt, rounded corners=1pt},
    etiq/.style={fill=udrojoosc, text=white, font=\fontsize{6.8pt}{8pt}\selectfont\bfseries, inner sep=2pt, anchor=north west},
    nota/.style={draw=udplata, fill=udfondo, font=\fontsize{7pt}{8.4pt}\selectfont, align=left, inner sep=3pt}]

  % --- participantes (x en cm) ---
  \def\xU{0.0}   \def\xC{2.95}  \def\xT{5.9}   \def\xD{8.85}  \def\xM{11.8}  \def\xA{14.75}
  \node[part]      (U) at (\xU,0) {\textbf{Usuario}\\(estudiante)};
  \node[parthecho] (C) at (\xC,0) {Cuaderno\\(Jupyter / Colab)};
  \node[part]      (T) at (\xT,0) {\texttt{torchvision}\\(datos)};
  \node[part]      (D) at (\xD,0) {\texttt{DataLoader}\\(lotes)};
  \node[part]      (M) at (\xM,0) {\texttt{Modelo}\\(\texttt{nn.Module})};
  \node[part]      (A) at (\xA,0) {\texttt{Adam}\\(optimizador)};
  \foreach \p/\x in {U/\xU,C/\xC,T/\xT,D/\xD,M/\xM,A/\xA}
    \draw[vida] (\x,-0.35) -- (\x,-16.05);

  % --- marcos de repetición (se dibujan primero para quedar al fondo) ---
  \draw[marco] (\xC-0.45,-5.75) rectangle (\xA+0.55,-13.15);
  \node[etiq] at (\xC-0.45,-5.75) {loop \;\;para cada época (3)};
  \draw[marco, dashed] (\xC-0.25,-6.3) rectangle (\xA+0.35,-11.70);
  \node[etiq] at (\xC-0.25,-6.3) {loop \;\;para cada lote (64 imágenes)};

  % --- mensajes (etiqueta sobre la flecha, sin tapar la línea) ---
  \tikzset{et/.style={txt, above=1.6pt}}
  \draw[sync]    (\xU,-1.5)  -- node[et]{ejecutar celdas} (\xC,-1.5);
  \draw[sync]    (\xC,-2.2)  -- node[et]{\texttt{cargar\_dataset("MNIST")}} (\xT,-2.2);
  \draw[retorno] (\xT,-2.8)  -- node[et]{tensores $N{\times}1{\times}28{\times}28$} (\xC,-2.8);
  \draw[sync]    (\xC,-3.4)  -- (\xC+0.55,-3.4) -- (\xC+0.55,-3.85) -- (\xC,-3.85);
  \node[txt, anchor=west] at (\xC+0.65,-3.62) {\texttt{preparar\_datos()}: 55\,000 / 5\,000 / 10\,000; $x'=(x-\mu)/\sigma$};
  \draw[sync]    (\xC,-4.55) -- node[et]{\texttt{TensorDataset}, lotes de 64} (\xD,-4.55);
  \draw[sync]    (\xC,-5.2)  -- node[et]{\texttt{Adam(parámetros, lr=1e-3)}} (\xA,-5.2);

  \draw[retorno] (\xD,-7.40)  -- node[et]{\texttt{xb, yb}} (\xC,-7.40);
  \draw[sync]    (\xC,-8.00)  -- node[et]{\texttt{zero\_grad()}} (\xA,-8.00);
  \draw[sync]    (\xC,-8.60)  -- node[et]{\texttt{modelo(xb)}} (\xM,-8.60);
  \draw[retorno] (\xM,-9.15) -- node[et]{\emph{logits}} (\xC,-9.15);
  \draw[sync]    (\xC,-9.70)  -- (\xC+0.55,-9.70) -- (\xC+0.55,-10.05) -- (\xC,-10.05);
  \node[txt, anchor=west] at (\xC+0.65,-9.87) {\texttt{F.cross\_entropy(logits, yb)} $\to$ pérdida};
  \draw[sync]    (\xC,-10.80) -- node[et]{\texttt{pérdida.backward()}} (\xM,-10.80);
  \draw[sync]    (\xC,-11.35) -- node[et]{\texttt{opt.step()}} (\xA,-11.35);

  \draw[sync]    (\xC,-12.25) -- node[et]{\texttt{evaluar(validación)}} (\xM,-12.25);
  \draw[retorno] (\xM,-12.80) -- node[et]{pérdida y exactitud} (\xC,-12.80);

  \draw[sync]    (\xC,-13.80) -- node[et]{\texttt{evaluar(prueba)}, \emph{una sola vez}} (\xM,-13.80);
  \draw[retorno] (\xM,-14.35) -- node[et]{matriz de confusión; pares} (\xC,-14.35);
  \draw[sync]    (\xC,-14.95) -- node[et]{\texttt{torch.save} y recargar} (\xM,-14.95);
  \draw[retorno] (\xC,-15.60) -- node[et]{\texttt{[OK]} / \texttt{[REVISAR]}} (\xU,-15.60);

  \node[nota, anchor=north east] at (\xA+0.55,-16.50) {Paso 5: repite el bucle con $\eta\in\{10^{-1},10^{-3},10^{-5}\}$, 2 épocas cada una.};
\end{tikzpicture}

% dentro de un entorno figure. Usa solo la paleta de ud-informe (udsable, udplata, udrojo, udazul).
\begin{tikzpicture}[yscale=0.78, font=\fontsize{7.6pt}{9pt}\selectfont,
    part/.style={draw=udsable!70, fill=white, rounded corners=2pt, minimum width=2.25cm, minimum height=0.7cm, align=center, inner sep=2pt},
    parthecho/.style={part, draw=udrojo, fill=udrojo!8, font=\fontsize{7.6pt}{9pt}\selectfont\bfseries},
    vida/.style={udplata!80, dash pattern=on 2.5pt off 2.5pt, line width=0.5pt},
    sync/.style={-{Latex[length=1.8mm]}, line width=0.65pt, color=udsable},
    retorno/.style={-{Latex[length=1.8mm, open]}, line width=0.55pt, color=udsable!65, dash pattern=on 3pt off 2pt},
    txt/.style={fill=white, inner sep=1.2pt, text=udsable},
    marco/.style={draw=udrojo!80, line width=0.7pt, rounded corners=1pt},
    etiq/.style={fill=udrojoosc, text=white, font=\fontsize{6.8pt}{8pt}\selectfont\bfseries, inner sep=2pt, anchor=north west},
    nota/.style={draw=udplata, fill=udfondo, font=\fontsize{7pt}{8.4pt}\selectfont, align=left, inner sep=3pt}]

  % --- participantes (x en cm) ---
  \def\xU{0.0}   \def\xC{2.95}  \def\xT{5.9}   \def\xD{8.85}  \def\xM{11.8}  \def\xA{14.75}
  \node[part]      (U) at (\xU,0) {\textbf{Usuario}\\(estudiante)};
  \node[parthecho] (C) at (\xC,0) {Cuaderno\\(Jupyter / Colab)};
  \node[part]      (T) at (\xT,0) {\texttt{torchvision}\\(datos)};
  \node[part]      (D) at (\xD,0) {\texttt{DataLoader}\\(lotes)};
  \node[part]      (M) at (\xM,0) {\texttt{Modelo}\\(\texttt{nn.Module})};
  \node[part]      (A) at (\xA,0) {\texttt{Adam}\\(optimizador)};
  \foreach \p/\x in {U/\xU,C/\xC,T/\xT,D/\xD,M/\xM,A/\xA}
    \draw[vida] (\x,-0.35) -- (\x,-16.05);

  % --- marcos de repetición (se dibujan primero para quedar al fondo) ---
  \draw[marco] (\xC-0.45,-5.75) rectangle (\xA+0.55,-13.15);
  \node[etiq] at (\xC-0.45,-5.75) {loop \;\;para cada época (3)};
  \draw[marco, dashed] (\xC-0.25,-6.3) rectangle (\xA+0.35,-11.70);
  \node[etiq] at (\xC-0.25,-6.3) {loop \;\;para cada lote (64 imágenes)};

  % --- mensajes (etiqueta sobre la flecha, sin tapar la línea) ---
  \tikzset{et/.style={txt, above=1.6pt}}
  \draw[sync]    (\xU,-1.5)  -- node[et]{ejecutar celdas} (\xC,-1.5);
  \draw[sync]    (\xC,-2.2)  -- node[et]{\texttt{cargar\_dataset("MNIST")}} (\xT,-2.2);
  \draw[retorno] (\xT,-2.8)  -- node[et]{tensores $N{\times}1{\times}28{\times}28$} (\xC,-2.8);
  \draw[sync]    (\xC,-3.4)  -- (\xC+0.55,-3.4) -- (\xC+0.55,-3.85) -- (\xC,-3.85);
  \node[txt, anchor=west] at (\xC+0.65,-3.62) {\texttt{preparar\_datos()}: 55\,000 / 5\,000 / 10\,000; $x'=(x-\mu)/\sigma$};
  \draw[sync]    (\xC,-4.55) -- node[et]{\texttt{TensorDataset}, lotes de 64} (\xD,-4.55);
  \draw[sync]    (\xC,-5.2)  -- node[et]{\texttt{Adam(parámetros, lr=1e-3)}} (\xA,-5.2);

  \draw[retorno] (\xD,-7.40)  -- node[et]{\texttt{xb, yb}} (\xC,-7.40);
  \draw[sync]    (\xC,-8.00)  -- node[et]{\texttt{zero\_grad()}} (\xA,-8.00);
  \draw[sync]    (\xC,-8.60)  -- node[et]{\texttt{modelo(xb)}} (\xM,-8.60);
  \draw[retorno] (\xM,-9.15) -- node[et]{\emph{logits}} (\xC,-9.15);
  \draw[sync]    (\xC,-9.70)  -- (\xC+0.55,-9.70) -- (\xC+0.55,-10.05) -- (\xC,-10.05);
  \node[txt, anchor=west] at (\xC+0.65,-9.87) {\texttt{F.cross\_entropy(logits, yb)} $\to$ pérdida};
  \draw[sync]    (\xC,-10.80) -- node[et]{\texttt{pérdida.backward()}} (\xM,-10.80);
  \draw[sync]    (\xC,-11.35) -- node[et]{\texttt{opt.step()}} (\xA,-11.35);

  \draw[sync]    (\xC,-12.25) -- node[et]{\texttt{evaluar(validación)}} (\xM,-12.25);
  \draw[retorno] (\xM,-12.80) -- node[et]{pérdida y exactitud} (\xC,-12.80);

  \draw[sync]    (\xC,-13.80) -- node[et]{\texttt{evaluar(prueba)}, \emph{una sola vez}} (\xM,-13.80);
  \draw[retorno] (\xM,-14.35) -- node[et]{matriz de confusión; pares} (\xC,-14.35);
  \draw[sync]    (\xC,-14.95) -- node[et]{\texttt{torch.save} y recargar} (\xM,-14.95);
  \draw[retorno] (\xC,-15.60) -- node[et]{\texttt{[OK]} / \texttt{[REVISAR]}} (\xU,-15.60);

  \node[nota, anchor=north east] at (\xA+0.55,-16.50) {Paso 5: repite el bucle con $\eta\in\{10^{-1},10^{-3},10^{-5}\}$, 2 épocas cada una.};
\end{tikzpicture}

\caption{Diagrama de secuencia del cuaderno resuelto. Línea continua: llamada; línea discontinua: retorno.}\label{fig:secuencia}
\end{figure}

%% ---------------------------------------------------------------------------------------------
\section{Bloque a bloque}

\begin{table}[H]
\centering\small
\caption{Qué hace cada bloque del script y qué idea enseña.}
\begin{tabularx}{\linewidth}{@{}l L L@{}}
\toprule
Bloque & Qué hace & Idea que enseña \\
\midrule
Paso 0 & Importa librerías, fija la semilla 42, elige CPU o GPU y define \texttt{verificar()} & Reproducibilidad y puntos de control \texttt{[OK]} / \texttt{[REVISAR]} \\
Paso 1 & Carga el dataset y separa 55\,000 / 5\,000 / 10\,000 imágenes; \textbf{normaliza con la media y la desviación del entrenamiento}; si no hay red, usa \texttt{digits} de scikit-learn & No filtrar información del futuro hacia el entrenamiento; el cuaderno funciona sin internet \\
Paso 2 & MLP: aplana la imagen $\to$ 128 neuronas ReLU $\to$ 10 \textit{logits}; entropía cruzada y Adam & Una línea base antes de la CNN; $784\cdot128+128+128\cdot10+10=101\,770$ parámetros \\
Paso 3 & CNN: dos bloques Conv $3\times3$ + ReLU + \textit{max-pooling} (16 y 32 filtros), \textit{dropout} y capa lineal; sin relleno, $28\to26\to13\to11\to5$ & Pesos compartidos: 12\,810 parámetros, ocho veces menos que el MLP \\
Paso 4 & Evalúa en prueba una sola vez; matriz de confusión y los tres pares más confundidos & Estimación honesta y análisis de errores \\
Paso 5 & Entrena con tres tasas de aprendizaje y 2 épocas cada una & Efecto de $\eta$: muy alta no aprende, muy baja es lenta \\
Paso 6 & Guarda el \texttt{state\_dict}, lo recarga y compara la exactitud & Persistir el modelo sin reentrenar \\
\bottomrule
\end{tabularx}
\end{table}

%% ---------------------------------------------------------------------------------------------
\section{Resultados que dejó el script}

\begin{table}[H]
\centering\small
\caption{Salidas guardadas en el cuaderno del profesor (3 épocas, Adam, $\eta=10^{-3}$, lotes de 64).}
\begin{tabularx}{\linewidth}{@{}l r r r L@{}}
\toprule
Modelo & Parámetros & Validación & Prueba & Observación \\
\midrule
MLP 784-128-10 & 101\,770 & 96,86\,\% & no evaluado & línea base \\
CNN (2 convoluciones) & 12\,810 & 98,58\,\% & 98,67\,\% & 133 errores de 10\,000; par más frecuente 4$\to$9 (12 veces) \\
\bottomrule
\end{tabularx}

\vspace{0.8em}
\begin{tabular}{@{}l r r@{}}
\toprule
Tasa de aprendizaje $\eta$ (2 épocas) & Validación & Pérdida de validación \\
\midrule
$10^{-1}$ & 10,42\,\% & 2,305 \\
$10^{-3}$ & 98,00\,\% & 0,058 \\
$10^{-5}$ & 81,22\,\% & 1,183 \\
\bottomrule
\end{tabular}
\end{table}

%% ---------------------------------------------------------------------------------------------
\section{Valoración del script}

\begin{table}[H]
\centering\small
\caption{Aciertos y límites encontrados al analizarlo.}
\begin{tabularx}{\linewidth}{@{}l L@{}}
\toprule
Tipo & Observación \\
\midrule
Acierto & La prueba se usa \textbf{una sola vez} y la normalización no usa información de validación ni de prueba. \\
Acierto & Hay una línea base (MLP) antes de la CNN y puntos de control automáticos en cada bloque. \\
Acierto & Reconoce que un resultado distinto del esperado debe anotarse y explicarse, no ocultarse. \\
Límite & \textbf{Las comparaciones no son pareadas.} \texttt{preparar\_datos} crea el \texttt{DataLoader} con un único generador compartido; su estado avanza en cada corrida, así que
          ``la misma semilla'' no repite el orden de los lotes. Con una corrida por configuración, una diferencia pequeña puede ser ruido. \\
Límite & \textbf{El Paso 5 es poco informativo:} con $\eta=10^{-1}$ queda en $\approx10$\,\% (azar), pero solo se probaron 3 valores, 2 épocas y una semilla, y la CNN no lleva
          \textit{BatchNorm}: no se sabe si el fallo es de la tasa o de la falta de normalización. \\
Límite & \textbf{Solo 3 épocas:} la validación aún sube (98,58\,\% en la última), así que son cifras de un modelo sin converger. El MLP no se evalúa en prueba; la comparación MLP--CNN es solo en validación. \\
Límite & Los rangos de exactitud del texto son orientativos: la nota inicial admite que se escribieron sin poder descargar MNIST y hay que contrastarlos con la propia ejecución. \\
\bottomrule
\end{tabularx}
\end{table}

Estos límites son los que se tuvieron en cuenta en el desarrollo propio (punto 4): cargador de datos por corrida, varias semillas, un barrido de $\eta$ más amplio con un control sin \textit{BatchNorm}
y evaluación en prueba de todas las redes comparadas.

\end{document}
